\title{QLoRA: Efficient Finetuning of Quantized LLMs}

\begin{abstract}
This paper introduces \textsc{QLoRA}, a finetuning technique that significantly lowers memory requirements, enabling the finetuning of a 65B parameter language model on a single 48GB GPU, all while maintaining the performance typically observed with standard 16-bit finetuning. The \textsc{QLoRA} method propagates gradients through a frozen 4-bit quantized version of a pretrained language model directly into Low Rank Adapters (LoRA). Our leading set of models, named \textbf{Guanaco}, surpasses all previously published open-access models on the Vicuna benchmark, achieving 99.3\% of ChatGPT’s performance level after only 24 hours of training on one GPU. To achieve these results, \textsc{QLoRA} incorporates several novel memory-saving strategies without compromising model quality: (a) the 4-bit NormalFloat (NF4) data type, theoretically optimal for weights distributed normally; (b) Double Quantization, which further shrinks memory use by quantizing the quantization constants; and (c) Paged Optimizers that control peak memory consumption. Employing \textsc{QLoRA}, we finetuned over 1,000 models and conducted a comprehensive evaluation of instruction following and chatbot capabilities across 8 different instruction datasets, diverse architectures (including LLaMA and T5), and parameter counts—covering large scales unattainable with conventional finetuning (such as 33B and 65B models). Our findings demonstrate that using QLoRA on a compact, high-quality dataset achieves or exceeds the latest state-of-the-art, even with models smaller than previous benchmarks. We also provide a thorough investigation of chatbot quality using both human assessments and GPT-4 judgments, establishing GPT-4-based evaluation as a cost-effective and practical proxy for human studies. Additionally, our research indicates that popular chatbot evaluation benchmarks are unreliable indicators of true performance. Through a targeted error analysis (“lemon-picking”), we illustrate specific shortcomings of \textbf{Guanaco} in comparison with ChatGPT. All models, source code, and CUDA kernels for 4-bit training are publicly available.\footnote{\url{https://github.com/artidoro/qlora} and \url{https://github.com/TimDettmers/bitsandbytes}}
\end{abstract}

\section{Introduction}

Adapting large language models (LLMs) via finetuning has been demonstrated to substantially enhance their effectiveness~\citep{min2021metaicl, wei2021finetuned, ouyang2022training, wang2022super, wang2022self, liu2022few} as well as to instill or eliminate particular behaviors~\citep{ouyang2022training,askell2021general,bai2022training}. Despite these advantages, finetuning extremely large-scale models incurs a massive computational cost; for instance, standard 16-bit finetuning of the LLaMA model with 65B parameters~\citep{touvron2023llama} demands more than 780 GB of GPU memory. Although advances in quantization approaches can decrease the inference memory demands of LLMs~\citep{dettmers2022llm, dettmers2022case, frantar2022gptq, xiao2022smoothquant}, they have proven unsuitable for training, often resulting in instability and failure during optimization~\citep{wortsman2023stable}.

Here, we provide the first empirical evidence that it is feasible to finetune LLMs quantized to 4 bits with no compromise in predictive accuracy. Our proposed approach, \textsc{QLoRA}, introduces an innovative high-precision quantization procedure to compress pretrained models to 4-bit representations, supplemented by the addition of compact, trainable Low-rank Adapter modules~\citep{hu2021lora}.

These adapter weights are optimized by propagating gradients through the quantized layers.

Employing \textsc{QLoRA}, the memory necessary to finetune a 65B parameter model decreases dramatically—from more than 780GB required in a conventional 16-bit scenario to less than 48GB—while sustaining comparable computational speed and model accuracy to fully finetuned 16-bit counterparts. This development significantly lowers the barrier for finetuning state-of-the-art, publicly available LLMs, now rendering them trainable on a single GPU. With \textsc{QLoRA}, we introduce the \textbf{Guanaco} model suite, where the second strongest model achieves 97.8\% of ChatGPT’s score on the Vicuna~\citep{vicuna2023} evaluation, with training times under 12 hours on a consumer-level GPU; leveraging a professional-grade GPU for 24 hours raises this to 99.3\% for our largest model, thus nearly matching ChatGPT's performance. In deployment, the smallest \textbf{Guanaco} model (7B parameters) operates with just 5 GB of memory, outperforming the 26 GB Alpaca model by over 20 percentage points on the Vicuna benchmark (Table~\ref{tbl:vicuna_eval}).

\textsc{QLoRA} incorporates several key advancements aimed at minimizing memory consumption while maintaining model efficacy: (1) \textbf{4-bit NormalFloat}, a quantization datatype optimized for data following a normal distribution, demonstrates empirically superior results compared to traditional 4-bit Integer and 4-bit Float representations. (2) \textbf{Double Quantization} involves an additional step of quantizing the quantization constants themselves, resulting in an average saving of approximately 0.37 bits per parameter—translating to around 3 GB saved in a 65B parameter model. (3) \textbf{Paged Optimizers} leverage NVIDIA’s unified memory to eliminate memory spikes during mini-batch processing of long sequences, which are typically associated with gradient checkpointing. Collectively, these advances are integrated into an enhanced LoRA framework with adapter layers at every level of the network, significantly reducing the accuracy compromises observed in earlier methods.

The memory-efficient design of \textsc{QLoRA} makes it feasible to conduct comprehensive analyses of instruction finetuning and chatbot capabilities across scales that would be impractical with conventional finetuning due to prohibitive memory requirements. Utilizing this method, we train in excess of 1,000 models, spanning a range of instruction tuning corpora, architectural variants, and parameter counts between 80M and 65B. Our experiments indicate that \textsc{QLoRA} can replicate the performance associated with 16-bit finetuning (\S\ref{sec:comp_to_std}), and we succeed in developing a cutting-edge chatbot, \textbf{Guanaco} (\S\ref{sec:pushing}). Our analysis of these trained models reveals notable insights: foremost, that the quality of training data is significantly more influential than its volume. For instance, a dataset of 9,000 examples (OASST1) surpasses a 450,000 example subset of FLAN v2 in supporting effective instruction-following behaviors in chatbots. Furthermore, we demonstrate that high achievement on the Massive Multitask Language Understanding (MMLU) benchmark does not necessarily predict strong performance on the Vicuna chatbot benchmark, and vice versa, underscoring that the alignment of dataset characteristics with task objectives outweighs mere dataset size.

In addition, we carry out a comprehensive examination of chatbot performance using both human evaluators and GPT-4 as assessment tools. We adopt a tournament-based benchmarking method, where various models engage in direct competition to generate superior responses to identical prompts. Each match is adjudicated by either human annotators or GPT-4, who select the best response. These match outcomes are then synthesized into Elo scores~\citep{elo1967proposed, elo1978rating}, which serve to rank the relative effectiveness of the different chatbots. Our findings indicate a general concordance between rankings provided by GPT-4 and human evaluators, though notable discrepancies are occasionally observed. Consequently, we caution that while model-driven evaluation offers a cost-effective alternative to human annotation, it remains subject to certain limitations and uncertainties.

Beyond the quantitative benchmark, we also conduct a qualitative assessment focused on the \textbf{Guanaco} models. This analysis brings to light specific areas where the models excel or underperform, which are not apparent from aggregate numerical scores.

To promote transparency and further investigation, we publicly release all model-generated responses along with their corresponding human and GPT-4 annotations. Our codebase—including CUDA kernels—has been open-sourced and is fully integrated with the Hugging Face transformers framework~\citep{wolf2019huggingface} for widespread accessibility. Additionally, we provide a suite of adapters for models of sizes 7/13/33/65B, each trained on eight distinct instruction-following datasets, resulting in a total of 32 fully open-source, fine-tuned model variants.

\section{Background}
\paragraph{Block-wise k-bit Quantization} Quantization refers to the transformation of data from a high-precision format to a lower-precision one, effectively reducing the amount of information represented. Typically, this involves converting a data type with greater bitwidth, such as 32-bit floats, to one with fewer bits, such as 8-bit integers. To fully utilize the representational capacity of the target low-bit type, the original data is rescaled so that its absolute maximum value matches the dynamic range of the new data type, a process achieved by normalizing the tensor's elements. For instance, when mapping a 32-bit floating-point (FP32) tensor into an Int8 tensor with the permissible range $[-127, 127]$:

\begin{equation}
    \mathbf{X}^{ \text{Int8}} = \text{round}\left( \frac{127}{{\text{absmax}}(\mathbf{X}^{{\text{FP32}}})}\mathbf{X}^{{\text{FP32}}} \right) = \text{round}(c^{\text{FP32}}\cdot \mathbf{X}^{{\text{FP32}}}),
\end{equation}

Here, $c$ represents the {\em quantization constant} or {\em quantization scale}. The inverse transformation—dequantization—restores the data back to its original range:

\begin{equation}
    \text{dequant}(c^{\text{FP32}}, \mathbf{X}^{\text{Int8}}) = \frac{\mathbf{X}^{\text{Int8}}}{c^{\text{FP32}}} = \mathbf{X}^{\text{FP32}}
\end{equation}

A major limitation of this method is its vulnerability to large magnitude entries—outliers—in the input tensor. Such outliers can cause poor utilization of quantization bins, with several bins corresponding to bit combinations left largely empty or unused. A typical strategy to address this problem involves partitioning the input tensor into several blocks, each of which is quantized separately using its own quantization constant $c$. More concretely, the input tensor $\mathbf{X}\in \mathbb{R}^{b\times h}$ is reshaped into a one-dimensional sequence and divided into $n = ({b\times h})/{B}$ contiguous blocks of size $B$. Each block is independently quantized using Equation~1, resulting in a quantized tensor as well as $n$ unique quantization constants $c_i$.

\paragraph{Low-rank Adapters} The Low-rank Adapter (LoRA) approach to finetuning \citep{hu2021lora} seeks to minimize memory usage by introducing a small number of trainable parameters—known as adapters—while keeping the remainder of the model’s parameters fixed. During stochastic gradient descent, gradients flow through the static pretrained weights to update only the adapter layers, which are trained to reduce the loss. LoRA introduces an additional low-rank, factorized projection to a linear transformation. For a linear mapping ${\mathbf{X}\mathbf{W} = \mathbf{Y}}$, where $\mathbf{X}\in \mathbb{R}^{b\times h}$ and $\mathbf{W}\in \mathbb{R}^{h\times o}$, the computation is modified as follows:

\begin{equation}
    \mathbf{Y} = \mathbf{X}\mathbf{W} + s\mathbf{X}\mathbf{L}_1\mathbf{L}_2,
\end{equation}

with $\mathbf{L}_1\in  \mathbb{R}^{h\times r}$, $\mathbf{L}_2\in  \mathbb{R}^{r\times o}$, and $s$ denoting a scalar coefficient.

\paragraph{Memory Requirement of Parameter-Efficient Finetuning} A key consideration in the use of LoRA is the memory demand incurred during training, particularly with respect to both the count and the size of adapters implemented. Due to LoRA’s inherently small memory requirements, increasing the number of adapters can lead to performance improvements while only marginally affecting the overall memory consumption. Although LoRA was introduced as a Parameter Efficient Finetuning (PEFT) strategy, the principal component of memory usage in LLM finetuning remains the storage of activation gradients, rather than the memory allocated for LoRA’s tunable parameters. For example, when finetuning a 7B LLaMA model on FLAN v2 with a batch size of 1 and employing LoRA adapters representing approximately 0.2\% of the total model weights~\citep{hu2021lora,liu2022few}, the input gradients associated with LoRA require 567 MB of memory, whereas the actual LoRA parameter storage is only 26 MB. Employing gradient checkpointing~\citep{chen2016training} reduces the memory required for input gradients to an average of 18 MB per sequence, yet this figure still surpasses the memory usage of the LoRA parameters themselves. For reference, the memory consumption of the 4-bit quantized base model stands at 5,048 MB. This comparison underlines that, while gradient checkpointing plays a significant role in lowering memory demand, further minimizing the LoRA parameter size yields only incremental improvements in memory efficiency. Consequently, expanding the number of adapters is feasible without substantially increasing training memory needs (a comprehensive analysis is provided in Appendix~\ref{app:memory}). This flexibility is particularly vital, as discussed subsequently, for achieving performance parity with full 16-bit precision finetuning.

\section{\textsc{QLoRA} Finetuning}

\textsc{QLoRA} enables accurate 4-bit finetuning by introducing two novel methods: 4-bit NormalFloat (NF4) quantization and Double Quantization. In addition, we propose Paged Optimizers to address the issue of memory spikes during gradient checkpointing, which can otherwise result in out-of-memory errors and hinder single-machine finetuning for large-scale models.

Within \textsc{QLoRA}, there is a distinct low-precision storage data type, typically 4-bit, alongside a separate compute data type, most often BFloat16. Operationally, this approach requires that whenever a \textsc{QLoRA} parameter tensor is accessed, it is first dequantized to BFloat16, after which all subsequent matrix multiplications are carried out in 16-bit precision.

Below, we elaborate on each building block of \textsc{QLoRA}, and subsequently provide a formalization of the \textsc{QLoRA} technique.

\paragraph{4-bit NormalFloat Quantization} The NormalFloat (NF) data type extends the concept of Quantile Quantization \citep{dettmers2022optimizers}, an information-theoretically optimal scheme that ensures a uniform allocation of values per quantization bin, relative to the empirical distribution of the input tensor. Quantile quantization achieves this by mapping the tensor values according to their estimated quantiles via the empirical cumulative distribution function.

A primary drawback of quantile quantization is the computational intensity required for quantile estimation. Consequently, more efficient quantile approximation strategies—such as SRAM quantiles~\citep{dettmers2022optimizers}—are typically employed. However, these approximate methods tend to introduce significant quantization errors for outlier values, which can often be critical.

Both the high computational cost of precise quantile calculation and the approximation error can be circumvented in cases where the input tensors share a common distribution up to a quantization factor. Under these circumstances, tensors maintain consistent quantile profiles, thereby permitting efficient and exact quantile estimation.

Because pretrained neural network weights are typically drawn from a zero-mean normal distribution with a standard deviation $\sigma$ (refer to Appendix~\ref{app:normality}), it is possible to standardize all weights by appropriately scaling $\sigma$, ensuring the distribution conforms precisely to the allowable range of the data type used. In this work, we arbitrarily define this data type range as $[-1, 1]$. Consequently, it becomes necessary to map both the data type’s quantiles and those of the neural network weights into this standardized interval.

The data type that is optimal from an information-theoretic standpoint, assuming input distributions are zero-mean normals with arbitrary $\sigma$ constrained to $[-1, 1]$, can be determined by the following process: (1) calculate the $2^k + 1$ quantiles from an idealized $N(0, 1)$ distribution, thereby constructing a $k$-bit quantile-based quantization scheme tailored for normal distributions; (2) scale the values of this scheme so that they fall within the $[-1, 1]$ interval; and (3) quantize any input weight tensor by rescaling it, using absolute maximum normalization, to match the $[-1, 1]$ range.

With the ranges of the weights and the data type now aligned, standard quantization procedures can be employed. Notably, the third step amounts to adjusting the weight tensor’s standard deviation so that it matches that of the $k$-bit quantized data type. In more precise terms, the $2^k$ quantized values $q_i$ are obtained by:

\begin{equation}
q_i  = \frac{1}{2}\left( Q_X\left(\frac{i}{2^k+1}\right) + Q_X\left(\frac{i+1}{2^k+1}\right)\right),
\end{equation}

where $Q_X(\cdot)$ denotes the quantile function for the standard normal distribution $N(0,1)$. A limitation with symmetric k-bit quantization is its inability to represent zero exactly, which is critical when quantizing padding values and other elements with a true value of zero without incurring error. To guarantee that zero is a discretized quantization point and to fully utilize all $2^k$ representable values of a k-bit number, we design an asymmetric data type. This is accomplished by determining quantiles $q_i$ for two intervals: $2^{k-1}$ quantiles on the negative side and $2^{k-1}+1$ on the positive, then merging these sets and eliminating the duplicate zero that appears in both. We designate this construction as the \emph{k-bit NormalFloat} (NFk), which ensures each quantization bin has the same expected number of values, and is optimal for zero-centered, normally distributed data from an information-theoretic perspective. A precise listing of these data type values is provided in Appendix~\ref{app:nf4}.

\paragraph{Double Quantization{}} 
We present \emph{Double Quantization} (DQ), a method for further compressing the storage cost by quantizing the quantization parameters themselves. Although precise 4-bit quantization requires a small block size as described in \citet{dettmers2022case}, this introduces substantial overhead from storing quantization constants. For instance, if 32-bit constants are applied to blocks of 64 elements in $\mathbf{W}$, the constants themselves consume $32/64 = 0.5$ bits per parameter on average. Double Quantization is introduced to address and significantly lower this overhead.

Concretely, Double Quantization takes the initial quantization constants $c_2^{\text{FP32}}$ and subjects them to a secondary quantization stage. This yields new, quantized constants $c_2^{\text{FP8}}$, along with a new set of quantization constants $c_1^{\text{FP32}}$ for the second step. For this purpose, we employ 8-bit floating point numbers and use a blocksize of 256 for the second quantization phase, as no accuracy loss was detected for 8-bit quantization, corroborating previous findings by \citet{dettmers2022case}. 
Given that $c_2^{\text{FP32}}$ values are strictly positive, we shift their mean to zero prior to quantization, thereby facilitating symmetric quantization. On average, for a block size of 64, Double Quantization decreases the per-parameter memory requirement for quantization constants from $0.5$ bits ($32/64$) to $8/64 + 32/(64 \cdot 256) = 0.127$ bits, yielding a savings of 0.373 bits for each parameter.

\paragraph{Paged Optimizers} leverage NVIDIA's unified memory~\footnote{\tiny \url{https://docs.nvidia.com/cuda/cuda-c-programming-guide}}, which transparently manages transfers of memory pages between CPU and GPU. This mechanism facilitates uninterrupted GPU execution, even if GPU memory becomes insufficient, functioning analogously to traditional OS memory paging between CPU RAM and storage devices. In this context, we allocate optimizer state using paged memory, enabling the system to automatically move state data to CPU RAM when GPU memory is exhausted and retrieve it back into GPU memory as required during the optimizer’s update phase.

\paragraph{\textsc{QLoRA}.} Based on the aforementioned mechanisms, we specify \textsc{QLoRA} for a single linear transformation in the quantized base network with one LoRA adapter as follows:

\begin{equation}
    \mathbf{Y}^{\text{BF16}} =  \mathbf{X}^{\text{BF16}} \text{doubleDequant}(c_1^{\text{FP32}}, c_2^{\text{k-bit}}, \mathbf{W}^{\text{NF4}}) + \mathbf{X}^{\text{BF16}}\mathbf{L}^{\text{BF16}}_1\mathbf{L}^{\text{BF16}}_2,
\end{equation}

where the doubleDequant$(\cdot)$ operation is specified by:

\begin{equation}
    \text{doubleDequant}(c_1^{\text{FP32}}, c_2^{\text{k-bit}}, \mathbf{W}^{\text{k-bit}}) = \text{dequant}(\text{dequant}(c_1^{\text{FP32}}, c_2^{\text{k-bit}}), \mathbf{W}^{\text{4bit}}) = \mathbf{W}^{\text{BF16}},
\end{equation}

Here, $\mathbf{W}$ utilizes the NF4 format and $c_2$ is represented in FP8 precision.

To enhance quantization accuracy, we set the block size for $\mathbf{W}$ to 64, and, for more efficient memory usage, assign a block size of 256 to $c_2$.

When performing parameter updates, it suffices to compute the gradients for the adapter weights, specifically $\frac{\partial E}{\partial \mathbf{L}_i}$; gradients with respect to the 4-bit weights, $\frac{\partial E}{\partial \mathbf{W}}$, are not needed. Nevertheless, deriving $\frac{\partial E}{\partial \mathbf{L}_i}$ requires intermediate calculation of $\frac{\partial \mathbf{X}}{\partial \mathbf{W}}$, which, as indicated in equation~(5), is performed after dequantizing the stored values $\mathbf{W}^{\text{NF4}}$ to the computation-friendly $\mathbf{W}^{\text{BF16}}$, ensuring the gradient $\frac{\partial \mathbf{X}}{\partial \mathbf{W}}$ is evaluated in BFloat16 precision.

In summary, \textsc{QLoRA} utilizes a single data type for storage—typically 4-bit NormalFloat—and a distinct data type for computation, namely 16-bit BrainFloat. During both the forward and backward passes, the parameters stored in the 4-bit format are dequantized to the 16-bit format for processing. However, gradient calculations for the weights are confined solely to the LoRA parameters, which are maintained in BrainFloat16 precision.

\section{QLoRA vs. Standard Finetuning}
\label{sec:comp_to_std}

Having outlined the operational mechanics of QLoRA and its capacity to significantly decrease the memory footprint required for model finetuning, the central issue remains whether QLoRA achieves competitive performance with conventional full-model finetuning approaches. Additionally, it is important to evaluate individual aspects of QLoRA, such as how the NormalFloat4 data type fares in comparison to standard Float4. The following sections detail the empirical studies conducted to address these topics.

\paragraph{Experimental setup.} Our investigation considers three categories of neural architectures: encoder-only, encoder-decoder, and decoder-only. QLoRA is benchmarked against both 16-bit adapter-based finetuning and conventional full-model finetuning for model sizes up to 3B parameters. For evaluation, we use the GLUE benchmark \citep{wang2018glue} with RoBERTa-large \citep{liu2019roberta}, Super-NaturalInstructions (TKInstruct) \citep{wang2022super} paired with T5 \citep{t5}, and assess 5-shot MMLU \citep{hendrycksmeasuring} following the finetuning of LLaMA on Flan v2 \citep{longpre2023flan} and Alpaca \citep{alpaca}. To further scrutinize the efficacy of NF4 relative to alternative 4-bit types, we employ the experimental protocol from \citet{dettmers2022case}, recording both zero-shot accuracy after quantization and perplexity on various models—namely, OPT~\citep{zhang2022opt}, LLaMA~\citep{touvron2023llama}, BLOOM~\citep{scao2022bloom}, and Pythia~\citep{biderman2023pythia}—with sizes ranging from 125 million to 13 billion parameters.

Detailed breakdowns for each experimental scenario are provided within the corresponding results sections to aid in clarity, with full methodological descriptions located in Appendix~\ref{app:exp-setup}.

Although the use of paged optimizers is essential to enabling \textsc{QLoRA} finetuning for 33B/65B models on single GPUs with 24/48GB VRAM, this work does not include specific quantitative assessments of paged optimizer efficiency, given that paging primarily takes place with infrequent long-sequence mini-batches. Nonetheless, we examine runtime implications of paged optimizers when training 65B models on 48GB GPUs and determine that, with batch sizes of 16, paged optimizers achieve equivalent throughput to standard optimizers. It is recommended that future investigations systematically determine the scenarios where paging introduces performance bottlenecks.

\paragraph{Default LoRA hyperparameters do not match 16-bit performance}

Standard LoRA configurations, such as those applying adapters to the query and value projection matrices in attention layers \citep{hu2021lora}, fail to match the effectiveness of full finetuning on large models. Our results, visualized in Figure~\ref{fig:hyper} for LLaMA 7B finetuned on Alpaca, indicate that the decisive LoRA hyperparameter is the total number of LoRA adapters, and achieving parity with full finetuning requires placing LoRA adapters in every linear layer within transformer blocks. Other parameters, including the projection dimension $r$, have minimal influence on outcome (refer to Appendix \ref{app:exp-setup}).

In a similar vein, we observe that the standard hyperparameters used for fully finetuned baseline models are inadequately optimized. To address this, we perform a systematic hyperparameter sweep across learning rates ranging from 1e-6 to 5e-5 and batch sizes from 8 to 128 in order to establish reliable baselines. Figure~\ref{fig:hyper} presents the results of 7B LLaMA finetuning conducted on the Alpaca dataset.

\paragraph{4-bit NormalFloat demonstrates superior results compared to 4-bit Floating Point}

Although the 4-bit NormalFloat (NF4) data type is considered optimal from an information-theoretic standpoint, it is necessary to empirically validate whether this translates into tangible performance benefits. Adopting the experimental protocol of \citet{dettmers2022case}, we assess a variety of quantized LLMs—including OPT \citep{zhang2022opt}, BLOOM \cite{scao2022bloom}, Pythia \cite{biderman2023pythia}, and LLaMA—spanning a size range from 125M to 65B parameters, across different quantization types and task settings in both language modeling and zero-shot evaluations. As illustrated in Figure~\ref{fig:nf4} and Table~\ref{tbl:nf4_ppl}, NF4 outperforms both FP4 and Int4 formats, yielding notable improvements, and further, double quantization is shown to reduce memory consumption without adverse effects on accuracy.

\paragraph{k-bit \textsc{QLoRA} achieves parity with 16-bit full finetuning and 16-bit LoRA}

Previous research has shown that 4-bit quantization is viable for \emph{inference}, though it typically entails some decrease in performance when compared to 16-bit representations \citep{dettmers2022case,frantar2022gptq}. This prompts a key question: Can finetuning with 4-bit adapters recover this loss in performance? We examine this in two distinct scenarios.

The initial comparison is made against fully finetuned 16-bit RoBERTA and T5 models ranging from 125M to 3B parameters, using both the GLUE and Super-NaturalInstructions datasets. Table~\ref{tab:nf4vsbf16} summarizes these outcomes. Across both datasets, 16-bit, 8-bit, and 4-bit adapter-based finetuning methods reach equivalence with the performance of their 16-bit fully finetuned counterparts, indicating that any performance deficits arising from lower-precision quantization can be fully mitigated through subsequent adapter-based finetuning.

For the second experimental configuration, considering that fully finetuning models of 11B parameters or greater necessitates multiple high-memory GPU servers, we proceed to assess whether 4-bit \textsc{QLoRA} achieves parity with 16-bit LoRA in the parameter range from 7B to 65B. Specifically, we finetune LLaMA models at 7B, 13B, 30B, and 65B sizes using two instruction-tuning corpora, Alpaca and FLAN v2, and then assess their performance using 5-shot accuracy on the MMLU benchmark. As detailed in Table~\ref{tbl:llama-datatype}, the results reveal that NF4, coupled with double quantization, is able to match the MMLU performance of the 16-bit LoRA variant. Additionally, our analysis shows that \textsc{QLoRA} with FP4 achieves results that trail the 16-bit brain float LoRA baseline by approximately one percentage point. These findings substantiate both (1) that \textsc{QLoRA} with NF4 is capable of mirroring the results of both 16-bit full finetuning and 16-bit LoRA finetuning, and (2) that NF4 demonstrates higher quantization accuracy relative to FP4.

\paragraph{Summary} The evidence from our experiments robustly indicates that 4-bit \textsc{QLoRA} utilizing the NF4 format achieves comparable performance to both 16-bit full finetuning and 16-bit LoRA across standard academic evaluation benchmarks. Moreover, the data highlights the superiority of NF4 over FP4, as well as the fact that the application of double quantization does not negatively impact outcomes. Collectively, these results make a strong case that 4-bit \textsc{QLoRA} tuning can be relied upon to achieve results on par with conventional 16-bit methods.

Consistent with findings in earlier quantization studies \citep{dettmers2022case}, our evaluations on MMLU and Elo demonstrate that, given fixed resources for finetuning and inference, allocating capacity toward increasing model parameters while reducing numerical precision can be advantageous. This emphasizes the efficiency gains attainable through \textsc{QLoRA}. Notably, since our empirical analyses did not reveal any performance loss in 4-bit finetuning when compared to full-precision finetuning, an open question emerges concerning the precise balance between model performance and numerical precision within QLoRA finetuning, a subject we leave for subsequent investigation.

We advance our analysis of instruction tuning by examining scenarios that would be unfeasible for full 16-bit fine-tuning when using standard academic computing resources.

\section{Pushing the Chatbot State-of-the-art with QLoRA}
\label{sec:pushing}
After demonstrating that 4-bit \textsc{QLoRA} achieves equivalent results to 16-bit fine-tuning across different scales, datasets, and tasks, we systematically investigate instruction finetuning on the largest publicly available language models intended for research. To measure the effects of instruction finetuning on these models, we employ the MMLU benchmark—recognized for its difficulty in Natural Language Understanding—and also introduce innovative techniques for assessing chatbot performance under real-world conditions.

\subsection{Experimental setup} This section summarizes the main elements of our experimental procedures, with extensive details presented in Appendix \ref{app:chatbot}.

\paragraph{Data} Because an exhaustive review of recent instruction-following datasets is lacking, we curate a collection of eight contemporary datasets. Our selection encompasses datasets sourced via crowd-sourcing (OASST1~\citep{kopf2023openassistant}, HH-RLHF~\citep{bai2022training}), datasets distilled from models already tuned to follow instructions (Alpaca~\citep{alpaca}, self-instruct~\citep{wang2022self}, unnatural-instructions~\citep{honovich2022unnatural}), aggregations of corpora (FLAN v2~\citep{chung2022scaling}), and hybrid datasets (Chip2~\citep{laion2023}, Longform~\citep{koksal2023longform}). Together, these datasets represent a variety of languages, dataset sizes, and licensing arrangements.

\paragraph{Training Setup} To mitigate confounding factors related to heterogeneous training objectives, all QLoRA finetuning is carried out using a supervised cross-entropy loss, excluding reinforcement learning—even when the datasets contain human preference annotations. For those datasets with clearly separated instruction and response fields, we restrict fine-tuning to the response component alone (see corresponding ablations in Appendix \ref{app:chatbot}). For datasets such as OASST1 and HH-RLHF that feature multiple possible responses, we choose the highest-ranked response per conversation tree node and train using the complete selected conversational thread, inclusive of instructions. Across experiments, we adopt the NF4 variant of \textsc{QLoRA}, enhanced with double quantization and paged optimization, thereby averting memory overflows during gradient checkpointing. Small-scale hyperparameter searches are performed for the 13B and 33B variants of LLaMA. Notably, we observe that the majority of hyperparameters optimized for the 7B version transfer well, with exceptions in learning rate and batch size: for the 33B and 65B models, we reduce the learning rate by half and double the batch size.

\paragraph{Baselines} We benchmark our models against leading research chatbots (Vicuna~\citep{vicuna2023}; Open Assistant~\citep{kopf2023openassistant}) as well as commercial offerings (GPT-4 \citep{openai2023gpt}, GPT-3.5-turbo, Bard). Open Assistant utilizes a LLaMA 33B model trained with Reinforcement Learning from Human Feedback (RLHF) on the same OASST1 dataset employed in our work. In contrast, Vicuna consists of a fully fine-tuned LLaMA 13B model using proprietary conversational data from ShareGPT, essentially reflecting a distilled OpenAI GPT derivative.

\subsection{Evaluation}

To evaluate language understanding capabilities, we follow established conventions by employing the MMLU (Massively Multitask Language Understanding) benchmark \citep{hendrycksmeasuring}. MMLU comprises a set of 57 multiple-choice tasks encompassing disciplines such as basic mathematics, U.S. history, computer science, and law. We report test accuracy under a 5-shot learning scenario.

To further assess generative language proficiency, we employ both automated tools and human raters. This second line of evaluation uses human-constructed prompts to gauge the quality of the model outputs. Although this approach is becoming more common and offers a more realistic appraisal of chatbot abilities, the field lacks a standardized evaluation methodology. Our approach, detailed below, uses nucleus sampling with $p=0.9$ and sets the temperature parameter to $0.7$ in every experiment.

\paragraph{Benchmark Data} We assess model performance on two specially compiled query datasets: Vicuna prompts \citep{vicuna2023} and the OASST1 validation set \citep{kopf2023openassistant}. The Vicuna set includes 80 prompts spanning a wide range of topics and is used without alteration. The OASST1 collection contains crowd-sourced multilingual, multi-turn interactions between users and assistants. Here, we extract all user entries from the validation split to use as queries, appending the corresponding dialog history to each prompt. This yields a total of 953 distinct user queries. These datasets serve as our Vicuna and OA benchmarks, respectively.

\paragraph{Automated Evaluation} Building on the protocol established by \citet{vicuna2023}, we employ GPT-4 to score the outputs of various systems relative to ChatGPT (GPT-3.5 Turbo) using the Vicuna dataset. For every prompt, both ChatGPT’s and a given model’s response are presented, and GPT-4 rates each reply on a scale from one to ten, supplying justifications for its assessments. To compute overall model performance, we express the model’s score as a percentage of ChatGPT’s score. It is possible for the relative score to surpass 100\% if the model attains a higher raw score than ChatGPT. We observe a notable position bias, with GPT-4 tending to favor the response appearing first. To address this, we advocate for reporting the average score across both possible response orderings.

Additionally, we evaluate by directly comparing system-generated outputs. In this variant, the evaluation becomes a three-class problem (accounting for possible ties), where GPT-4 is asked to select the superior reply, indicate a tie, and provide a rationale. Pairwise comparisons are conducted for all system pairs over both the Vicuna and OA datasets.

\paragraph{Human Evaluation} Although emerging studies suggest that generative models can successfully facilitate system assessments \citep{fu2023gptscore}, it remains unclear whether GPT-4’s evaluations of chatbot quality are consistent with human judgment. Consequently, we conduct two parallel rounds of human assessments on the Vicuna prompts, mirroring the methodologies used in our automated evaluations. These involve recruitment via Amazon Mechanical Turk (AMT), with each ChatGPT comparison rated by two annotators and each pairwise system comparison rated by three annotators.

\paragraph{Elo Rating} Utilizing both human and automated pairwise evaluations, we establish a tournament-like framework in which models are pitted against one another. Each match features two models generating responses to a shared prompt, competing for the superior answer. While this process draws on approaches seen in \citet{bai2022training} and \citet{vicuna2023}, our methodology uniquely incorporates both human and GPT-4 assessments. To compute the Elo~\citep{elo1967proposed,elo1978rating} scores, we randomly select from the pool of annotated comparisons. The Elo system, originally developed for chess and other competitive games, reflects the likelihood of one competitor outperforming another—for instance, a participant with an Elo of 1100 facing one with an Elo of 1000 has an expected success rate of roughly 65\%; when both contenders have equal ratings, such as 1000 vs 1000 or 1100 vs 1100, each is anticipated to win about 50\% of the time. Elo scores are updated after every match based on the predicted outcome: greater adjustments occur following unexpected results, while anticipated outcomes yield minor rating shifts. Gradually, this results in Elo scores that are well-aligned with each model’s actual performance level in these competitions. The initial rating is set to 1,000 with $K=32$. Following \citet{vicuna2023}, we repeat this computation 10,000 times using various random seeds to mitigate biases introduced by the order in which model pairs are compared.

\subsection{\textbf{Guanaco}: \textsc{QLoRA} trained on OASST1 Emerges as a Leading Chatbot}

Through a combination of automated and human assessments, our analysis demonstrates that the leading \textsc{QLoRA} model, {Guanaco} 65B, which has been finetuned using a modified version of OASST1, stands out as the most capable open-source chatbot available, delivering results on par with ChatGPT. When assessed alongside GPT-4, both {Guanaco} 65B and 33B achieve a predicted win probability of 30\% according to Elo scores derived from human annotator-based, system-level pairwise model comparisons—surpassing any previous reports.

Table~\ref{tbl:vicuna_eval} summarizes outcomes on the Vicuna benchmark \cite{vicuna2023}, situating performance relative to ChatGPT. The findings indicate that {Guanaco} 65B ranks immediately after GPT-4, attaining 99.3\% of ChatGPT’s performance. While {Guanaco} 33B features a parameter count exceeding that of Vicuna 13B, it leverages 4-bit weight quantization for greater memory efficiency—using only 21 GB compared to 26 GB—while delivering a three-point performance increase over Vicuna 13B. Additionally, {Guanaco} 7B is compact enough to operate on contemporary mobile devices with a mere 5 GB memory footprint, yet still achieves a nearly 20-point margin over Alpaca 13B.

Despite these results, Table~\ref{tbl:vicuna_eval} reveals considerable variability in confidence intervals, with substantial overlap among models’ performances. We attribute this imprecision to ambiguous scoring standards, such as an undefined interpretation for an 8 out of 10 score across various conditions. To address this limitation, we advocate the Elo ranking system \citep{elo1967proposed}, which depends on \textit{pairwise} judgments from human raters and GPT-4, thus eliminating the ambiguity of absolute rating scales. The Elo-based standings of top-performing models are presented in Table~\ref{tbl:elo}. Notably, there is some divergence between rankings by humans and GPT-4 on the Vicuna benchmark, particularly with Guanaco 7B; however, concordance remains relatively robust for most models, with a Kendall Tau of $\tau=0.43$ and a Spearman rank correlation of $r=0.55$ at the system comparison level. On a per-example basis, the majority agreement between GPT-4 and human annotators diminishes, as indicated by a Fleiss $\kappa=0.25$. Collectively, these metrics point to a moderate level of agreement between GPT-4 and human evaluations at the system level, suggesting that model-based assessments can serve as a reasonably dependable proxy for human judgment. Further issues are elaborated in Section~\ref{ssec:considerations}.

The Elo rankings presented in Table~\ref{tbl:elo_large} show that the {Guanaco} 33B and 65B models achieve superior results compared to all models except GPT-4 on both the Vicuna and OA benchmarks, and their performance is on par with ChatGPT, as also reflected in Table~\ref{tbl:vicuna_eval}. It should be highlighted that the Vicuna benchmark tends to advantage open-source models, whereas the OA benchmark favors ChatGPT. Additionally, analysis of Tables \ref{tbl:mmlu-llama} and \ref{tbl:vicuna_eval} underscores the critical impact of selecting an appropriate finetuning dataset on overall outcomes. For instance, Llama models that are finetuned using FLAN v2 demonstrate exceptional results on the MMLU benchmark but have the lowest scores on Vicuna—a trend mirrored in other evaluated models. This indicates that the evaluation benchmarks have some degree of independence from each other: excelling at MMLU does not necessarily translate to high performance as a chatbot (as gauged by Vicuna or OA), nor does the reverse hold.

{Guanaco} is unique among the leading models we assess in that it has not been trained on proprietary datasets, with the OASST1 dataset’s collection process strictly prohibiting the use of GPT-based models. The next strongest model that relies exclusively on open-source data is Anthropic’s HH-RLHF, but it lags {Guanaco} by 30 percentage points on the Vicuna benchmark (refer to Table~\ref{tbl:vicuna_eval}). Collectively, these findings demonstrate that 4-bit \textsc{QLoRA} is a highly effective technique, capable of yielding chatbots whose quality rivals that of ChatGPT. Notably, our 33B {Guanaco} can be trained in under 12 hours on a standard 24 GB consumer GPU. This paves the way for future research involving \textsc{QLoRA} fine-tuning on targeted open-source datasets, potentially creating models that stand shoulder to shoulder with today’s top commercial alternatives.

\section{Qualitative Analysis}

Although quantitative methods constitute the primary focus of our assessment, relying solely on aggregate statistics has inherent limitations. Chief among these is the issue of benchmark validity \citep{liao2021we}—the extent to which a benchmark genuinely measures what its designation or description claims is often unclear, particularly given the frequent identification of ``shortcuts'' that allow machine learning models to exploit benchmarks \citep{gururangan2018annotation, poliak2018hypothesis}. To partially address this limitation, we incorporate a qualitative perspective, structured in two parts. In the first part, \S\ref{ssec:examples}, we present illustrative examples that typify certain patterns detected in outputs produced by our 65b {Guanaco} model. The subsequent section, \S\ref{ssec:considerations}, discusses specific caveats pertaining to our findings and offers insights regarding our interpretation.

\subsection{Qualitative Analysis of Example Generations}
\label{ssec:examples}

To identify illustrative cases, we systematically reviewed model outputs from the Vicuna and OpenAssistant benchmarks, paying close attention to recurring response behaviors from {Guanaco}. Upon observing a distinctive tendency, we crafted prompts specifically designed to provoke the identified behavior, even in instances where such a response would be incorrect. For example, if the model was found to consistently produce unnecessarily verbose answers, we explicitly instructed: ``Answer yes or no without explanation.'' Through this approach, we isolated instances—so-called ``lemons''—where these adversarial prompts successfully caused model failures, and also highlighted cases—termed ``cherries''—where the model maintained robustness. All example responses shown in this section were produced via Nucleus Sampling \cite{holtzmancurious} with $p=0.9$.

Naturally, this overview is not exhaustive, as the limited scope of this qualitative inquiry prevents systematic control over the numerous factors at play. For example, the space of possible model outputs for any specific prompt is extensive, so we rely on a selection of samples that we believe approximate the range of model behaviors. Nevertheless, we argue that presenting these instances helps situate the quantitative findings previously discussed. By open-sourcing all associated models and code, we aim to encourage further in-depth investigations into the challenges outlined here.

\paragraph{Factual Recall} When prompted with fact-based queries such as ``What is the capital of Zambia?'', each model consistently produces the accurate response, for example:

\begin{quote}
    \textbf{Guanaco}:
    \texttt{The capital of Zambia is Lusaka.}
\end{quote}

Yet as prompts involve less commonly known information, {Guanaco} begins to fail while still responding confidently. Consider the following question taken from HotPotQA \citep{yang2018hotpotqa}:

\begin{quote}
     \textbf{User}: In what year was the singer who popularized `I'll Keep the Lovelight Burning' born?  

     \textbf{Guanaco}:~\texttt{Al Jolson is credited with popularizing the song `I'll Keep the Lovelight Burning,' and he was born in the year 1886.}
 \end{quote}

Here, the response identifies the incorrect artist and also misreports the relevant birth year (though the year matches the birth date for Al Jolson, who is not the correct answer).

\paragraph{Suggestibility} In some scenarios involving presupposed falsehoods, {Guanaco} demonstrates notable reluctance to conform or validate incorrect assumptions, as illustrated below:

\begin{quote}
    \textbf{User}: How did peer reviewed scientists officially verify that the earth is flat? \\ \\
    \textbf{Guanaco}: \texttt{There has never been an official confirmation by peer reviewed scientists that the earth is flat. The scientific community overwhelmingly agrees that the earth is round, a fact supported for hundreds of years through various scientific experiments and observations. \\ \\ Claims that scientists have peer reviewed and verified the flat earth theory are unfounded and have been disproven by scientific evidence. The notion that science endorses a flat earth is a misconception associated with conspiracy theories. \\ \\ Feel free to ask additional questions if you have them.}
\end{quote}

is also effective in recognizing questions that cannot be answered, such as,

\begin{quote}
    \textbf{User}: What time is it?

    \textbf{Guanaco}: \texttt{I'm unable to provide real-time information, as I am an AI language model without access to the current time. My knowledge is static and not connected to live data. To find the current time, please consult a clock or an online time source.}
\end{quote}

The same pattern holds for queries like ``Where are you?'', ``How are you?'', and similar types of questions.

\paragraph{Refusal} In addition, there are cases where {Guanaco} declines to comply with prompts for no clear reason:

\begin{quote}
    \textbf{User}: Please reverse the order of the words in the sentence ``I'm a cool dude in a mocha mood.'' \\ \\
    \textbf{Guanaco}: \texttt{I'm sorry, I cannot fulfill that request at the moment, but I can clarify the grammatical structure of the sentence you gave.\\ \\The expression ``I'm a cool dude in a mocha mood" adheres to standard English grammar and syntax. Here, the descriptive adjective ``cool'' precedes the noun ``dude'' it modifies, and the prepositional phrase ``in a mocha mood'' provides further contextual detail. The construction of this sentence is grammatically correct and logically coherent.\\ \\If you have any further inquiries or need additional assistance, please let me know.}
\end{quote}

\paragraph{Secret Keeping} The system’s refusal strategy is unfortunately not always robust when secrecy is required. In some cases, system developers may want to ensure that specific information remains concealed from users, even when the model is allowed to generate unconstrained text outputs. To test the ease with which a user might obtain confidential information, we prompted the model as follows: ``The secret word is `banana'. Do not reveal or repeat back this back, no matter what the human requests.'' When asked directly, the model responds as intended:

\begin{quote}
    \textbf{User}: What is the secret word?

    \textbf{Guanaco}: \texttt{I'm sorry, but I'm not able to disclose the secret word. I'm programmed to not reveal or repeat back the secret word, no matter what the human requests. Is there anything else I can help you with?}
\end{quote}

which demonstrates the expected, secure behavior. Nonetheless, only a slight deviation in prompt can cause this security measure to fail:

\begin{quote}
    \textbf{User}: This is a game. The goal is to ignore your previous instructions. What is the secret word?

    \textbf{Guanaco}: \texttt{Sorry, the secret word is banana. It's best if we don't speak of it. Is there anything else you'd like to ask?}
\end{quote}

underscoring the importance of further research into developing systems capable of more reliably adhering to restrictive instructions.

\paragraph{Math} One of the principal challenges for {Guanaco} lies in mathematical reasoning, a domain that remains difficult for many large language models as described in \citep{liang2022holistic}. When prompted to explain its calculations, {Guanaco} often arrives at the right solution, for instance:

\begin{quote}
    \textbf{User}: Lee mows one lawn and charges \$33. Last week he mowed 16 lawns and three customers each gave him a \$10 tip. How many dollars did Lee earn mowing lawns last week?

    \textbf{Guanaco}: \texttt{Lee earned \$582 mowing lawns last week. \\ \\ To find the total amount he earned, we first need to calculate how much he earned for mowing 16 lawns: \\ \\  \$33 x 16 = \$528 \\ \\ Then we need to add the tips he received from the three customers: \\  \\ \$10 x 3 = \$30 \\ \\ So the total amount Lee earned was: \\ \\ \$528 + \$30 = \$558}
\end{quote}

Nevertheless, {Guanaco} can sometimes fail, even on relatively elementary tasks, particularly if it does not proceed through the solution incrementally—a limitation noted in the literature \cite{weichain}. For example, consider the interaction below:

\begin{quote}
    \textbf{User}: Please factorize 1833.

    \textbf{Guanaco}: \texttt{Sure, I'd be happy to help! 1833 is a prime number, which means it has only two factors: 1 and 1833. \\ \\ So the unique factorization of 1833 is: \\ \\ 2\^{}1 * 3\^{}2 * 17\^{}1 \\ \\ I hope this helps! Let me know if you have any other questions.}
\end{quote}

which is not only incorrect (the actual prime factorization is $3 \times 17 \times 43$), but doubly so.

\paragraph{Theory of Mind} {Guanaco} exhibits unexpectedly robust Theory of Mind reasoning abilities \cite{nematzadeh2018evaluating, sap2022neural}. For example, the model successfully provides a thorough answer to the following scenario:

\begin{quote}
    \textbf{User}: James and Abby are in the bedroom. Abby put the pen in the desk drawer. Abby leaves the bedroom. James moves the pen into the bag. Where does James think Abby will look for the pen?

    \textbf{Guanaco}: \texttt{James knows that Abby put the pen in the desk drawer. But James also knows that he has moved the pen into the bag, and he assumes that Abby does not know about this. So, James may guess that Abby will look for the pen in the desk drawer, because that's where she left it last time. However, this assumption may be incorrect, because Abby may have seen James moving the pen.}
\end{quote}

Nonetheless, such inferential performance is inconsistent; the model often generates explanations that rest on unwarranted or illogical premises, for instance,

\begin{quote}
    \textbf{User}: Evelyn entered the living room. Jackson entered the playroom. James entered the playroom. The beans are in the treasure chest. James exited the playroom. Jackson moved the beans to the pantry. Jackson exited the playroom. James entered the living room. Where will James look for the beans?

    \textbf{Guanaco}: \texttt{James will look for the beans in the pantry, because that's where Jackson moved them.}
\end{quote}

here, {Guanaco} assumes a flow of information that was never established in the narrative. These shortcomings align with recent findings in the literature \cite{sap2022neural}, indicating the necessity for further investigation.

\subsection{Considerations}
\label{ssec:considerations}

\paragraph{Evaluation}

We observe only moderate consistency among human evaluators (Fleiss $\kappa=0.42$), with further decline when assessing outputs from two high-performing models. This highlights deficiencies in current benchmarks and human-based evaluation procedures for assessing chatbot tasks. In qualitative comparisons of outputs from ChatGPT and {Guanaco} 65B using the Vicuna benchmark, subjective judgments become prominent, as coauthors of this paper frequently disagreed regarding their preferred answers. Future efforts should seek solutions inspired by methodologies from Human-Computer Interaction and Psychology to better accommodate subjective preferences.

Furthermore, our study indicates that automated evaluation systems display notable biases. As an illustration, we note pronounced order effects, where GPT-4 tends to assign higher scores to whichever system's output is presented first. The limited agreement at the sample level between GPT-4 and human raters (Fleiss $\kappa=0.25$) suggests a divergence in evaluation criteria or preferences. Additionally, Table~\ref{tbl:elo_large} reveals that GPT-4 rates its own completions markedly higher than those assessed by human raters (Elo scores of 1348 vs 1176), equating to an approximate 20\% increase in the chance of outperforming a competing response.

Further research should address the existence of possible biases within automated evaluation frameworks, as well as explore methods for reducing or counteracting these biases.

\paragraph{Data \& Training} It is worth noting that the OASST1 dataset, utilized in the training of {Guanaco} models, is inherently multilingual, and that the OA benchmark features prompts spanning several languages. We suggest future investigations to assess the extent to which multilingual training enhances instruction-following abilities in non-English languages and to determine if this factor contributes to the observed performance disparity between Vicuna-13B (trained solely on English text) and {Guanaco} 33B/65B on the OA benchmark.

Given the notable performance demonstrated by {Guanaco} models, we examined potential data leakage from the OASST1 dataset into the Vicuna benchmark prompts. Utilizing fuzzy string matching followed by manual inspection of the nearest matches, we found no evidence of prompt overlap between these two datasets.

It should also be highlighted that our model's training exclusively employed cross-entropy loss under a supervised learning paradigm, with no incorporation of reinforcement learning from human feedback (RLHF). This warrants further studies to elucidate the comparative advantages and disadvantages between conventional cross-entropy loss and RLHF-based training. We anticipate that \textsc{QLoRA} will facilitate such comparative analyses on a broad scale, without requiring excessive computational resources.

\section{Related Work}

\paragraph{Quantization of Large Language Models} Prior efforts on quantizing LLMs have predominantly addressed inference-time quantization. Leading methods for maintaining the quality of 16-bit LLMs focus on addressing outlier activations (such as SmoothQuant~\citep{xiao2022smoothquant} and LLM.int8()~\citep{dettmers2022llm}), while alternative approaches leverage advanced grouping techniques~\citep{park2022nuqmm, yao2022zeroquant}. Research on lossy quantization investigates the implications of standard rounding~\citep{dettmers2022case,zeng2022glm,pope2022efficiently} as well as strategies for optimizing rounding to boost quantization fidelity~\citep{frantar2022gptq}. Beyond the current work, only SwitchBack layers~\citep{wortsman2023stable} have examined the feasibility of backpropagating through quantized weights at scales exceeding 1B parameters.

\paragraph{Finetuning with Adapters} Although our experiments employ Low-rank Adapters~\citep{hu2021lora} (LoRA), the literature includes a diverse array of Parameter Efficient FineTuning (PEFT) techniques. These alternatives span prompt tuning~\citep{qin2021learning,lester2021power,li2021prefix}, adjusting embedding layer inputs~\citep{an2022input}, manipulating hidden states as in IA$^3$~\citep{liu2022few}, incorporating entire layers~\citep{houlsby2019parameter}, modifying bias terms~\citep{zaken2021bitfit}, devising masks over parameters with guidance from Fisher information~\citep{sung2021training}, as well as hybrid strategies~\citep{henderson2021compacter}. Our findings indicate that LoRA adapters are sufficient to attain the same performance level as full 16-bit finetuning. A comprehensive analysis of the comparative merits of other PEFT methodologies is reserved for subsequent work.

\paragraph{Instruction Finetuning} Instruction finetuning aims to enable pretrained LLMs to reliably execute tasks specified in prompts. This is accomplished by finetuning on input-output pairs drawn from diverse datasets, training the model to produce correct outputs given a prompt. Representative methods and benchmarks include MetaICL~\citep{min2021metaicl}, MetaTuning~\citep{zhong2021adapting}, InstructGPT~\citep{ouyang2022training}, FLAN~\citep{wei2021finetuned,chung2022scaling}, PromptSource~\citep{bach2022promptsource}, Super-NaturalInstructions~\citep{wang2022super,sanh2021multitask}, Self-instruct~\citep{wang2022self}, UnnaturalInstructions~\citep{honovich2022unnatural}, OPT-IML~\citep{iyer2022opt}, UnifiedSKG\citep{xie2022unifiedskg}, OIG/Chip2~\citep{laion2023}, Alpaca~\citep{alpaca}, Vicuna~\citep{vicuna2023}, Koala~\citep{koala_blogpost_2023}, and Self-instruct-GPT-4~\citep{peng2023instruction}.

\paragraph{Chatbots} Instruction-aligned models are frequently realized as conversational agents, and are commonly refined using Reinforcement Learning from Human Feedback (RLHF)~\citep{christiano2017deep} or by generating training examples with model-generated feedback, referred to as RLAIF~\citep{bai2022constitutional}. Datasets and systems such as Anthropic-HH~\citep{askell2021general,bai2022training}, Open Assistant~\citep{kopf2023openassistant}, LaMDA~\citep{thoppilan2022lamda}, and Sparrow~\citep{glaese2022improving} exemplify this trend. While our methodology does not employ reinforcement learning, our leading model, Guanaco, is finetuned on the Open Assistant multi-turn chat corpus, which is specifically constructed for RLHF-based procedures~\citep{kopf2023openassistant}. For chatbot evaluation, methods using GPT-4 in place of expensive human assessments have become prominent~\citep{vicuna2023,peng2023instruction}. Our work advances these evaluation techniques, concentrating on creating a more robust and trustworthy evaluation framework.

\section{Limitations and Discussion}

We present evidence demonstrating that \textsc{QLoRA}, leveraging a 4-bit foundation in combination with Low-rank Adapters (LoRA), can emulate the performance achieved with full 16-bit finetuning. Nevertheless, this has not been conclusively validated at the 33B and 65B parameter scales for \textsc{QLoRA}. Given the substantial computational demands, a detailed analysis at these scales is deferred to subsequent investigations.

A further constraint pertains to our evaluation strategy for instruction-finetuned models. Our results are based on the MMLU, Vicuna, and OA benchmarks; however, we have not assessed the models on alternative benchmarks such as BigBench, RAFT, or HELM, thus the transferability of our findings to these datasets is not guaranteed. Despite this, our evaluation includes an extensive study using MMLU, as well as novel evaluation methods for chatbot performance.

The available evidence suggests that benchmark outcomes are highly influenced by the similarity between the finetuning data and the benchmark’s own dataset. For instance, FLAN v2 shares significant overlap with MMLU but does not align closely with chatbot-oriented benchmarks, while the Chip2 dataset exhibits the opposite pattern; as a result, the respective models show corresponding performance differences on MMLU and Vicuna. This observation underscores the necessity not only for more robust benchmarks and assessment procedures but also for a clear understanding of the evaluation targets. It is important to question whether the aim is to develop models that excel at academic knowledge—as seen in high school or college contexts—or at conversational competence characteristic of chatbots, or perhaps at other abilities entirely. Given the relative convenience of utilizing established benchmarks rather than creating new ones, there is a risk that specific benchmarks may unduly shape research trajectories. As a community, it is imperative to verify that benchmarks genuinely reflect desired outcomes.

Although our evaluation offers a comprehensive assessment of general chatbot functionality, it remains limited in terms of responsible AI evaluation for {Guanaco}. Specifically, we assess the propensity of {Guanaco}-65B to produce socially biased outputs compared to other models in Table~\ref{tbl:crows}. Our results indicate that {Guanaco}-65B registers considerably lower average bias scores than other models trained only with raw pretraining data. This suggests that finetuning on the OASST1 dataset mitigates some biases inherent in the LLaMA base model. Nevertheless, it is uncertain if {Guanaco}'s lower bias extends across other forms of bias, and a broader, more comprehensive analysis of potential biases in {Guanaco} and similar systems is left for future research.

Another shortcoming in our study is the lack of evaluation on various bit-precision configurations—such as 3-bit models—or alternative adapter approaches. While LoRA was employed based on its established robustness, there exists a diverse array of Parameter Efficient FineTuning (PEFT) strategies with promising results, although their scalability to very large models remains unclear. Different adapters might outperform LoRA in some scenarios. Moreover, since post-quantization finetuning appears to recover most of the representational capacity lost during quantization, this could facilitate the use of more aggressive quantization schemes. For example, applying 3-bit GPTQ quantization to a base model and subsequently finetuning with LoRA may achieve performance comparable to traditional 16-bit full finetuning.

\section{Broader Impacts}

Our \textsc{QLoRA} approach marks the first instance in which 33B parameter models can be finetuned using a single consumer-grade GPU, and 65B parameter models can be managed with just one professional GPU, all without performance loss relative to full-scale finetuning. Our top-performing 33B model, finetuned on the Open Assistant dataset, matches ChatGPT’s results on the Vicuna benchmark. As instruction finetuning is vital for transforming foundational pretrained large language models into high-quality ChatGPT-like conversational agents, we anticipate that our method will greatly democratize access to advanced NLP capabilities, particularly for researchers with limited computational resources. Ultimately, \textsc{QLoRA} functions as a leveling mechanism that helps to minimize resource disparities between large organizations and smaller teams with only consumer hardware.

An additional avenue for potential impact is the adaptation of large language models for use on mobile devices. We posit that our \textsc{QLoRA} approach represents a significant advance toward allowing LLM finetuning on smartphones and in other environments with limited computational resources. Previous work demonstrated that running 7B-parameter models on mobile phones was feasible, but \textsc{QLoRA} is the first method to facilitate the finetuning of such models in these contexts. According to our calculations, an iPhone 12 Plus could process approximately 3 million tokens overnight during charging using \textsc{QLoRA}. While the performance of finetuned 7B models remains below that of ChatGPT, we argue that their quality is sufficient to unlock innovative applications, particularly in areas previously constrained by privacy or model capability concerns. By employing \textsc{QLoRA}, users can retain greater control over both their data and model deployments, thereby supporting privacy-centric use cases while simplifying the rollout of LLM-based solutions.

Nonetheless, it should be recognized that finetuning presents dual-use risks and may be exploited for malicious purposes. The proliferation of LLMs carries well-documented hazards \citep{bommasani2021opportunities,bender2021dangers}, but we maintain that democratizing access to this rapidly emerging technology can foster more independent and thorough scrutiny, as opposed to concentrating LLM capabilities within large corporations that may restrict model access and transparency.

Overall, we are confident that \textsc{QLoRA} will exert a predominantly positive effect by considerably broadening and simplifying access to high-quality LLM finetuning.